\documentclass[11pt,letterpaper]{article}
\usepackage[margin=0.95in]{geometry}
\usepackage{amsmath,amssymb,booktabs,array,microtype}
\usepackage[T1]{fontenc}
\usepackage{lmodern}
\usepackage[colorlinks=true,linkcolor=blue,citecolor=blue,urlcolor=blue]{hyperref}
\setlength{\parindent}{0pt}\setlength{\parskip}{5pt}
\sloppy
\newcolumntype{L}[1]{>{\raggedright\arraybackslash}p{#1}}
\newcommand{\ms}{\,\mathrm{m\,s^{-2}}}

\title{Falsifiable Predictions of a de Sitter--MOND Clock Completion}
\author{Carl P.\ Zimmerman\\ \small Briar Creek Tech --- AI-assisted research draft; not peer reviewed}
\date{9 September 2026}

\begin{document}\maketitle

\begin{abstract}
We collect the sharply falsifiable predictions of a relativistic completion of Milgrom's dynamics in which
the modified-inertia sector is carried by a \emph{clock} --- a scalar entering through the clock-projected
derivative $Q=n^\mu\partial_\mu\varphi$ --- with the acceleration scale locked to the de Sitter/cosmological
scale, $a_0=c^2/2\pi L_{\rm dS}$, and cosmological dark matter supplied by the conserved Noether charge of
the clock-scalar's shift symmetry. Each prediction is stated as an observable, a threshold, and a
discriminator against dark matter or $\Lambda$CDM, and is separated into what is robust to the construction
and what is kernel-specific. Six are distinctive and currently testable: \textbf{(1)} the acceleration scale
is \emph{constant in cosmic time} (a flat $a_0(z)$), so the deep-MOND baryonic Tully--Fisher zero-point at
$z\simeq2.5$ is flat ($0.00$~dex) against a $+0.33$~dex $\Lambda$CDM expectation, and a robustly \emph{rising}
$a_0(z)$ (tracking the full $H(z)$, $\sim3\times$ by $z=2$) falsifies the theory; \textbf{(2)} an
\emph{External Field Effect}: internal velocity dispersions of dwarf spheroidals depend on Galactocentric
distance, rising $\sim\!1.9\times$ across $40$--$250$~kpc, where dark matter predicts \emph{no} dependence ---
the sharpest discriminator, and the theory reproduces the low dispersion of Crater~II; \textbf{(3)} in merging
clusters the collisionless charge-dust traces the galaxies, so the lensing mass separates from the gas (the
Bullet-cluster offset that pure MOND fails); \textbf{(4)} the tensor gravitational-wave speed equals $c$
exactly, and a \emph{subdominant, near-luminal scalar polarization} (amplitude set by the clock--metric
coupling) is predicted at a level consistent with current LIGO/pulsar-timing bounds; \textbf{(5)} the dark
matter is a conserved Noether charge and therefore \emph{cannot decay or annihilate}, forbidding indirect
detection lines; \textbf{(6)} the exact exponential kernel gives a parameter-free $1/r$ correction above the
flat rotation velocity, $v^2=\sqrt{G_bMa_0}+G_bM/4r+7(G_bM)^{3/2}/96r^2\sqrt{a_0}+\mathcal{O}(r^{-3})$. Each
number is produced by a committed, runnable script; every dimensional quantity is carried on both $a_0$
footings. \textbf{We do not claim the theory is established:} its dark sector carries an intrinsic
Big-Bang-Nucleosynthesis fine-tuning (companion paper), and its full relativistic health is a work in
progress; these are predictions of a \emph{candidate} class, offered so that data can confirm or kill it.
The construction was developed collaboratively in the associated public repository.
\end{abstract}

\section{The class}
The predictions below hold for relativistic MOND completions of the schematic form
$S=\int\sqrt{-g}[\tfrac{M^2}{2}R-\Lambda M^2-K(Q)+F(Q)\Theta+M^2a_0^2G(|V|/a_0)]+S_m$, with $n_\mu$ the unit
gradient of a clock field, $Q=n\!\cdot\!\partial\varphi$, $\Theta=\nabla\!\cdot\!n$, $G(y)=y^2+2(1+y)e^{-y}-2$
(so $G'(y)/2y=1-e^{-y}$), and matter minimally coupled to the single metric $g$. Two structural facts drive
the predictions: (i) $a_0$ is fixed by $\Lambda$, and (ii) the dark matter is the conserved charge
$a^3(-K_Q+3HF_Q)=C$, pressureless and clustering like cold dark matter. Predictions robust to these two facts
are marked $\star$; kernel-specific ones are marked $\dagger$.

\section{Predictions}

\textbf{P1 ($\star$) --- Flat $a_0(z)$; deep-MOND BTFR zero-point at high $z$.}
$a_0\propto\sqrt{\Lambda}$ is constant in cosmic time, so the deep-MOND baryonic Tully--Fisher relation
$V^4=G_bMa_0$ keeps a fixed zero-point with redshift. At $z\simeq2.5$ the predicted zero-point shift is
$0.00$~dex, versus a $+0.33$~dex expectation for a $\Lambda$CDM halo-abundance drift.
\emph{Falsifier:} a robust $a_0$ \emph{rising} with $z$ (tracking $H(z)$, $\sim\!3\times$ by $z=2$) kills the
Lock. \emph{Test:} a deep-MOND lensed rotator at $z\!\sim\!2$--$3$. [\texttt{L90}]

\textbf{P2 ($\star$) --- External Field Effect on dwarf spheroidals.}
The nonlinear MOND response violates the strong equivalence principle: a dwarf's internal dispersion is
suppressed by the external Galactic field. For a $10^6M_\odot$ dwarf, $\sigma$ rises $\sim\!1.9\times$ across
Galactocentric radius $40\to250$~kpc; Crater~II is pulled to $\sigma\simeq2.4$~km\,s$^{-1}$ (observed
$\simeq2.7$). \emph{Discriminator:} dark matter predicts \emph{no} $\sigma$--$R_{\rm gc}$ dependence. This is
the sharpest MOND-vs-dark-matter test. [\texttt{L89}]

\textbf{P3 ($\star$) --- Bullet-type mergers.}
The dark matter is a collisionless conserved charge, so in a cluster merger the lensing mass separates from
the X-ray gas and traces the galaxies --- the offset pure MOND famously fails. A centroid model with the
dust at $\sim\!6\times$ the baryons places the lensing peak on the galaxy side ($x_c/d\simeq0.15$--$0.36$ vs
$0.93$ for pure MOND-on-gas). \emph{Discriminator:} pure MOND fails; this class passes qualitatively.
[\texttt{L85}]

\textbf{P4 ($\star$) --- Gravitational-wave speed and polarization.}
$c_T=c$ exactly (the scalar/MOND terms contribute no graviton kinetic term), passing GW170817 structurally.
A subdominant, near-luminal scalar polarization (breathing/longitudinal, mass $\sim\!10^{-33}$~eV, amplitude
$\propto F_Q$) is predicted, consistent with current LIGO/Virgo/pulsar-timing polarization bounds.
\emph{Falsifier:} any measured $c_T\neq c$; \emph{discriminator:} detection of the scalar mode. [\texttt{L88}]

\textbf{P5 ($\star$) --- No dark-matter decay or annihilation.}
The dark matter is the Noether charge of an exact shift symmetry, hence exactly conserved. The class predicts
\emph{no} decay lines and \emph{no} annihilation signal, unlike a WIMP/particle relic.
\emph{Falsifier:} a confirmed dark-matter decay or annihilation detection. [\texttt{L81}]

\textbf{P6 ($\dagger$) --- A $1/r$ correction above the flat rotation velocity.}
Inverting the exact implicit law gives, for a point mass,
\[
v^2=\sqrt{G_bMa_0}+\frac{G_bM}{4r}+\frac{7(G_bM)^{3/2}}{96\,r^2\sqrt{a_0}}+\mathcal{O}(r^{-3}),
\]
a parameter-free correction with no fitting freedom. \emph{Discriminator:} a dark-matter halo gives a
profile-dependent (fittable) correction; this one is fixed by the kernel. \emph{Test:} precise rotation
curves in the transition region $g\sim a_0$. [\texttt{L94}]

\section{Consistency checks (must-match, not distinctive)}
The class must also reproduce, and does at the levels tested: the SPARC radial-acceleration relation
(exponential kernel, rms $0.16$~dex at zero per-galaxy parameters, ranking just below the $\nu_{\rm RAR}$
kernel; [\texttt{L92}]); the linear structure-growth amplitude ($\sigma_8=0.811$, matching $\Lambda$CDM to
machine precision because the charge-dust has $c_s^2=0$ and standard $G_{\rm eff}$; [\texttt{L93}]); galaxy
weak-lensing following the same relation as dynamics (no-slip $\Phi=\Psi$); and the Solar-System
post-Newtonian bounds, including the preferred-frame parameters that exclude the earlier AeST completion
($\alpha_1$ is structurally suppressed here rather than un-tunable; [\texttt{L91}]).

\section{Honest scope}
These are predictions of a candidate class, not of an established theory. Two costs are stated in the
companion work: the charge-dust carries an intrinsic $\sim\!24$-order Big-Bang-Nucleosynthesis fine-tuning
(the same charge sources an $a^{-6}$ stiff partner), and the full relativistic constraint algebra / health of
the propagating sector is still being derived. The wide-binary signal predicted by the exponential kernel
($\gamma_v\simeq1.03$--$1.06$) sits below the strongest claimed detections and is consistent with the
Newton-leaning analyses, so wide binaries chiefly test \emph{which} interpolating kernel is correct rather
than MOND versus dark matter. No observation is used here to assert the theory is true; the measured
quantities run only the other way, saying how large an effect a member of the class must show and what would
falsify it. Every number is produced by a committed script with checks that can fail
(\texttt{fable\_independent\_2026/L81,L85,L88--L94}); dimensional numbers are carried on both $a_0$ footings
($9.3619\times10^{-11}\ms$, $1.1279\times10^{-10}\ms$). AI-assisted research draft; not peer reviewed.

\begin{thebibliography}{9}
\bibitem{Milgrom1983} M.~Milgrom, ApJ \textbf{270}, 365 (1983).
\bibitem{BekensteinMilgrom1984} J.~Bekenstein and M.~Milgrom, ApJ \textbf{286}, 7 (1984).
\bibitem{SkordisZlosnik2021} C.~Skordis and T.~Z\l{}o\'snik, Phys.\ Rev.\ Lett.\ \textbf{127}, 161302 (2021).
\bibitem{McGaugh2016} S.~McGaugh, F.~Lelli and J.~Schombert, \emph{Radial acceleration relation in rotationally supported galaxies}, Phys.\ Rev.\ Lett.\ \textbf{117}, 201101 (2016).
\bibitem{Lelli2017} F.~Lelli, S.~McGaugh, J.~Schombert and M.~Pawlowski, \emph{One law to rule them all: the RAR in galaxies}, ApJ \textbf{836}, 152 (2017).
\bibitem{Clowe2006} D.~Clowe et al., \emph{A direct empirical proof of the existence of dark matter}, ApJ \textbf{648}, L109 (2006).
\bibitem{Chae2023} K.-H.~Chae, \emph{Breakdown of the Newton--Einstein standard gravity at low acceleration in wide binaries}, ApJ \textbf{952}, 128 (2023).
\bibitem{Abbott2017} LIGO/Virgo Collaborations, \emph{GW170817: Observation of gravitational waves from a binary neutron star inspiral}, Phys.\ Rev.\ Lett.\ \textbf{119}, 161101 (2017).
\bibitem{Zimmerman2026bbn} C.~P.~Zimmerman, \emph{An Intrinsic Big Bang Nucleosynthesis Fine-Tuning in Conserved-Charge Dark Matter from Relativistic MOND Clocks}, Zenodo (2026).
\end{thebibliography}

\end{document}
